\appendix
\counterwithin{figure}{chapter}

\chapter{Convertirea în formatul glb}

\begin{figure}[H]
    \begin{center}
        \includegraphics[height=18cm]{glb_export} \caption{Exportare în formatul glb din Blender}
        \label{fig:glb_export}
    \end{center}
\end{figure}


\chapter{Game Design Document}
\label{anex:gdd}
\noindent
\textbf{
    \Large{Eden - Short Game Design Document}}

\noindent
\large{Creat de: \emph{Cătălin Pleșu}}

\vspace{5mm}
\noindent This document is a quick overview of the game design decisions of the game \emph{Eden} (this is a short version of the usual Game Design Document).

\vspace{5mm}
\noindent\textbf{Story}

Actiunea are loc într-o pădure masivă.

\vspace{5mm}
\noindent\textbf {Characters}

Interaționăm cu lumea prin ochii unui caracter necunoscut la persoana I.

\vspace{5mm}
\noindent\textbf {Visual Arts}

Jocul va avea grafica 3D.

O mare parte din lume va fi generata procedural.

\begin{figure}[H]
    \begin{center}
        \includegraphics[width=10cm]{teren} \caption{Exemplu de teren generat procedural}
        \label{fig:teren}
    \end{center}
\end{figure}

\begin{figure}[H]
    \begin{center}
        \includegraphics[width=10cm]{plante} \caption{Exemplu de plante generate procedural}
        \label{fig:plante}
    \end{center}
\end{figure}

Imaginile folosite nu sunt generate de mine.

\vspace{5mm}
\noindent\textbf {Sound Effects (SFX) \& Music}

În joc nu va fi prezentă muzică, doar sunete care pot fi gasite în natură.

\vspace{5mm}
\noindent\textbf {Gameplay Features}

Jucătorul poate explora padurea, strângând diverse alimente pentru a nu muri de foame. Totodata va evita animalele periculoase care îl pot ataca.

\vspace{5mm}
\noindent\textbf {Level Design (LD)}

Joaca este cu lumea deschisă, jocul este castigat cand jucatorul gaseste un copac cu niste  fructe interzise

\vspace{5mm}
\noindent\textbf {User Interface (UI)}

Înafara de meniul de la începutul jocului nu va fi prezentă interfața grafică pentru a nu distrage jucătorul de la joc.

\vspace{5mm}
\noindent\textbf {Game Controls}

Control ca în majoritatea jocurilr cu taste 'wasd' și mouse.

\vspace{5mm}
\noindent\textbf {Target Platform / Minimum Hardware:}

Jocul va fi compilat pentru:
\begin{itemize}
    \item Linux;
    \item Windows.
\end{itemize}
Pentru o experienta placuta va fi nevoie de minimum:
\begin{itemize}
    \item 8 RAM
    \item 1 VRAM
    \item Intel i3
    \item 500 MB spatiu de stocare
\end{itemize}

\vspace{5mm}
\noindent\textbf {Team}

Cătălin Pleșu este singurul membru al echipei și va avea următoarele roluri:
\begin{itemize}
\item Designer;
\item Artist;
\item Programator;
\item Tester.
\end{itemize}

\vspace{5mm}
\noindent\textbf {Time}

Dezvoltarea jocului a început pe 21.02.2022 și a fost finisată pe 15.04.2022.
În total au fost 40 de zile lucrătoare, pentru a reuși trebuie să folosesc o tehnică
împărțite în 8 iterații utilizând metodologia Agile.

Considerând că majoritatea jocurilor dureza câțiva ani și au echipe de sute de
persoane nu pot fi foarte ambițios, si trebuie să am așteptări realiste.

\chapter{Capturi de ecran din joc}
\label{anex:gameplay}

\begin{figure}[H]
    \begin{center}
        \includegraphics[width=18cm]{gameplay-terrain} \caption{Vizualizarea terenului generat}
        \label{fig:gameplay-terrain}
    \end{center}
\end{figure}

\begin{figure}[H]
    \begin{center}
        \includegraphics[width=18cm]{gameplay-veiw} \caption{Parcurgerea hărții de către jucător}
        \label{fig:gameplay-veiw}
    \end{center}
\end{figure}

\begin{figure}[H]
    \begin{center}
        \includegraphics[width=18cm]{gameplay-fruits} \caption{Grămadă de fructe}
        \label{fig:gameplay-fruits}
    \end{center}
\end{figure}

\begin{figure}[H]
    \begin{center}
        \includegraphics[width=18cm]{gameplay-end} \caption{Jocul s-a finisat cu scorul 1632}
        \label{fig:gameplay-end}
    \end{center}
\end{figure}

